\section{Summary}
\subsubsection{End-to-end systems level comparisons under SWaP constraints}
Across the UAV DAA radar literature, feasibility has been demonstrated, but evidence remains fragmented across hardware prototypes, single-path processing demonstrations, and isolated GPU acceleration studies. As a result, there remains a gap in consolidated, systems-level research that illustrates which combinations of radar signal processing, monopulse angle estimation, and multi-target tracking algorithms are most effective when jointly optimised for sensitivity, latency, and throughput on GPU-based embedded platforms, and how these conclusions change with hardware configuration. The proposed evaluation of small-target sensitivity and end-to-end tracking on embedded GPU hardware directly supports our contributions on algorithm selection, low-level implementation trade-offs, and hardware-aware design recommendations. 

\subsubsection{Realistic small-target and encounter-focused evaluation}
The experimental design explicitly includes small moving aerial targets and encounter-like scenarios in both simulation and real data collection. Performance is measured not only by detection range but by false alarm behaviour, track continuity, and sensitivity in low-SNR regimes that better reflect small-aperture UAV radars. By extending prior single-path demonstrations into a comparative study of alternative algorithm combinations and fully onboard GPU implementations, this section motivates our contributions on algorithm selection and implementation choices, and on how these choices vary across hardware contexts.

\subsubsection{GPU-oriented onboard pipeline evidence, including data movement costs}
The thesis targets embedded GPU hardware and evaluates architectural choices that often dominate real-time performance, including host–device transfers, memory layout, streaming formats, kernel fusion opportunities, and synchronisation overhead. This thesis advances the literature by benchmarking low-level GPU implementations beyond range--Doppler and CFAR, quantifying end-to-end latency and throughput impacts, and assessing how these conclusions change across embedded hardware configurations.

\subsubsection{Hardware-aware design conclusions across embedded configurations}
Experiments are repeated across multiple hardware contexts (for example CPU-only, unified-memory embedded GPU, and discrete CPU+GPU settings where available) and across constrained power and clock modes. Testing these angle-estimation methods within a complete multi-target tracking pipeline, and measuring both sensitivity and track-quality impacts alongside compute cost, directly supports our contributions on algorithm selection, implementation trade-offs, and hardware-aware recommendations.